\chapter{Miara Wienera i twierdzenie Camerona--Martina}
\label{ch:cameron-martin}

\begin{uwaga}[O wersjach funkcjonału gaussowskiego]
Wzór Camerona--Martina zawiera $I(h')$, czyli element
$L^2(C_0([0,T]),\mu_T)$, a nie kanonicznie wyznaczoną
liczbę dla każdej trajektorii $w$. Wybór dowolnego
borelowskiego reprezentanta tej klasy daje tę samą
pochodną Radona--Nikodyma prawie wszędzie. W ogólności
nie istnieje ciągły funkcjonał liniowy na całej
przestrzeni $C_0([0,T])$, który realizowałby $I(h')$
dla dowolnego $h'\in L^2$.
\end{uwaga}

\section{Dlaczego miara na przestrzeni trajektorii jest innym obiektem niż całka Itô?}

Całka $\int_0^T H_t\,\dd W_t$ jest zmienną losową, a w przypadku
całkowania na zmiennym przedziale --- procesem. Natomiast dla funkcjonału
$F:C_0([0,T])\to\R$ zapis $\int F(w)\,\dd\mu_T(w)$ jest zwykłą całką
Lebesgue'a względem \emph{miary na przestrzeni funkcji}. Są to dwa
różne rodzaje całkowania. Pierwszy wykorzystuje przewidywalność i przyrosty
procesu, drugi jest uśrednieniem po wszystkich możliwych trajektoriach.

Na przykład $F(w)=\exp(-\int_0^T w(s)^2\,\dd s)$ jest funkcjonałem
trajektorii: dla każdej ciągłej $w$ całka po czasie istnieje klasycznie.
Wówczas $\int F\,\dd\mu_T$ jest liczbą z przedziału $(0,1]$.
Nie ma tu całki stochastycznej. W innym przykładzie
$F(w)=w(T)^2$ dostaniemy $\int F\,\dd\mu_T=T$.

\begin{definicja}[Kanoniczna przestrzeń Wienera]
\label{def:canonical-wiener}
Niech $T>0$ i
$C_0([0,T])=\{w\in C([0,T];\R):w(0)=0\}$ z normą
$\|w\|_\infty=\sup_{0\le t\le T}|w(t)|$ i
borelowskim $\sigma$-ciałem $\mathcal B_T$.
\emph{Miara Wienera} $\mu_T$ jest prawem standardowego procesu Wienera
na $C_0([0,T])$. Na przestrzeni kanonicznej proces współrzędnych
$B_t(w)=w(t)$ jest procesem Wienera pod miarą $\mu_T$.
\end{definicja}

W tomie II skonstruowaliśmy proces Wienera o ciągłych trajektoriach;
zatem odwzorowanie $\omega\mapsto(W_t(\omega))_{t\in[0,T]}$ jest
mierzalne jako odwzorowanie w $C_0([0,T])$ (wystarcza mierzalność
współrzędnych w chwilach wymiernych). Obraz miary probabilistycznej
przez to odwzorowanie jest $\mu_T$. Konstrukcja nie wymaga formalnej
„miary $\prod_t\dd w(t)$”.

\begin{twierdzenie}[Rozkłady skończeniewymiarowe i jednoznaczność miary]
\label{th:wiener-cyl}
Niech $0<t_1<\cdots<t_n\le T$, $x_0=0$, $t_0=0$, oraz
$p_t(z)=(2\pi t)^{-1/2}\exp(-z^2/(2t))$ dla $t>0$.
Dla każdego ograniczonego borelowskiego $g:\R^n\to\R$ zachodzi
\[
\int_{C_0([0,T])}g(w(t_1),\ldots,w(t_n))\,\dd\mu_T(w)
=\int_{\R^n}g(x_1,\ldots,x_n)
\prod_{j=1}^n p_{t_j-t_{j-1}}(x_j-x_{j-1})
\,\dd\lambda_n(x).
\]
Miara $\mu_T$ jest jedyną miarą probabilistyczną na
$C_0([0,T])$, dla której zachodzą te tożsamości.
\end{twierdzenie}
\begin{proof}
Dla procesu Wienera wektor przyrostów
$(W_{t_1},W_{t_2}-W_{t_1},\ldots,W_{t_n}-W_{t_{n-1}})$
ma niezależne współrzędne normalne z wariancjami
$t_1,t_2-t_1,\ldots,t_n-t_{n-1}$. Liniowa transformacja
$(y_1,\ldots,y_n)\mapsto(y_1,y_1+y_2,\ldots,y_1+\cdots+y_n)$
ma wyznacznik równy $1$. Wzór na gęstość wynika więc ze
zmiany zmiennych w $\R^n$.

Dla jednoznaczności wystarczy zauważyć, że funkcje
$e_q(w)=w(q)$, gdzie $q\in[0,T]\cap\Q$ (oraz $q=T$),
generują $\mathcal B_T$. Istotnie, dla $w,v\in C_0([0,T])$
z ciągłości wynika
$\|w-v\|_\infty=\sup_{q\in D}|w(q)-v(q)|$, gdzie
$D=([0,T]\cap\Q)\cup\{T\}$. Otwarte kule w normie
jednostajnej są zatem mierzalne względem
$\sigma(e_q:q\in D)$. Ponieważ $C_0([0,T])$ jest
separowalną przestrzenią metryczną, każde otwarte podzbiory
są przeliczalnymi sumami takich kul. Z drugiej strony
wszystkie $e_q$ są ciągłe. Stąd równość $\sigma$-ciał.
Zbiory cylindryczne w wymiernych chwilach tworzą
$\pi$-układ generujący $\mathcal B_T$; dwie miary
zgodne na nim są równe z twierdzenia o $\pi$- i $\lambda$-układach.
\end{proof}

\begin{przyklad}[Uśrednianie funkcjonału trajektorii]
Dla $F(w)=\int_0^T w(t)^2\,\dd t$ twierdzenie Tonellego daje
\[
\int F(w)\,\dd\mu_T(w)
=\int_0^T\E B_t^2\,\dd t
=\int_0^Tt\,\dd t=\frac{T^2}{2}.
\]
Wykorzystaliśmy tutaj \emph{dwie} miary: Lebesgue'a w czasie i
Wienera na przestrzeni trajektorii. Ich kolejność można zamienić,
ponieważ funkcja podcałkowa jest nieujemna.
\end{przyklad}

\section{Deterministyczna całka Wienera jako gaussowski funkcjonał trajektorii}

\begin{twierdzenie}[Izometria i funkcje charakterystyczne]
\label{th:pw-map}
Dla każdego $f\in L^2([0,T])$ istnieje zmienna losowa
$I(f)\in L^2(C_0([0,T]),\mu_T)$, jednoznaczna w $L^2$,
taka że dla funkcji schodkowej
$f=\sum_{j=0}^{n-1}c_j\ind{(t_j,t_{j+1}]}$ zachodzi
$I(f)(w)=\sum_j c_j(w(t_{j+1})-w(t_j))$.
Odwzorowanie $I:L^2([0,T])\to L^2(\mu_T)$ jest liniowe i
\[
\E_{\mu_T}[I(f)I(g)]=\int_0^T f(t)g(t)\,\dd t,
\qquad
\E_{\mu_T}e^{iuI(f)}=e^{-u^2\|f\|_2^2/2}.
\]
Dla dowolnych $f_1,\ldots,f_k\in L^2([0,T])$ wektor
$(I(f_1),\ldots,I(f_k))$ jest gaussowski.
\end{twierdzenie}
\begin{proof}
Dla schodkowych $f,g$ pierwsza tożsamość wynika z niezależności
przyrostów i z faktu, że kowariancja dwóch przyrostów ruchu Browna
jest równa długości przecięcia odpowiadających im przedziałów.
W szczególności $\|I(f)-I(g)\|_{L^2(\mu_T)}=\|f-g\|_2$.
Jeśli $f_n$ są schodkowe i zbiegają w $L^2$ do $f$,
to $I(f_n)$ jest ciągiem Cauchy'ego w $L^2(\mu_T)$;
jego granica nie zależy od wybranej aproksymacji.
Wzór na kowariancję przechodzi do granicy przez nierówność
Cauchy'ego--Schwarza. Każda kombinacja
$\sum_j u_jI(f_j)=I(\sum_j u_jf_j)$ jest granicą w $L^2$
zmiennych normalnych o wariancjach zbiegających do
$\|\sum_j u_jf_j\|_2^2$. Zbieżność funkcji charakterystycznych
i twierdzenie o ciągłości dowodzą normalności granicy.
To samo rozumowanie dla wszystkich liniowych kombinacji
daje wspólną gaussowskość wektora.
\end{proof}

\begin{uwaga}[Gdzie ukryta jest subtelność mierzalności?]
Dla każdej pojedynczej funkcji $f$ otrzymujemy klasę
równoważności $I(f)$, a nie jeden kanoniczny wzór punktowy,
który działałby jednocześnie dla wszystkich $f$ i wszystkich
trajektorii $w$. Nie możemy więc pisać bez wyjaśnienia
$I(f)(w)=\int f(t)w'(t)\,\dd t$: typowa trajektoria
Wienera nie posiada pochodnej. W dowodach poniżej każda
występująca całka stochastyczna jest zdefiniowana jako
granica w $L^2$, a wzory dotyczą równości prawie na pewno.
\end{uwaga}

\section{Przestrzeń Camerona--Martina: które przesunięcia są dopuszczalne?}

\begin{definicja}[Przestrzeń Camerona--Martina]
\label{def:CMspace}
Niech
\[
\calH_T=\left\{h\in C_0([0,T]):
\begin{array}{l}
 h\text{ jest absolutnie ciągła oraz}\\[-2pt]
 h'\in L^2([0,T])
\end{array}\right\}.
\]
Wprowadzamy iloczyn skalarny
$\langle h,k\rangle_{\calH_T}=\int_0^T h'(s)k'(s)\,\dd s$
oraz normę $\|h\|_{\calH_T}=\|h'\|_2$.
\end{definicja}

\begin{twierdzenie}[Struktura Hilberta i jądro reprodukujące]
\label{th:cm-hilbert}
Przestrzeń $\calH_T$ jest przestrzenią Hilberta,
a odwzorowanie $D:h\mapsto h'$ jest izometrią
$\calH_T$ na $L^2([0,T])$. Ponadto
\[
|h(t)-h(s)|\le\sqrt{|t-s|}\,\|h\|_{\calH_T}
\quad(s,t\in[0,T]),
\]
a funkcjonał wartości w chwili $t$ jest reprezentowany
przez $k_t(s)=\min(s,t)$:
$h(t)=\langle h,k_t\rangle_{\calH_T}$.
\end{twierdzenie}
\begin{proof}
Z definicji absolutnej ciągłości i podstawowego twierdzenia
rachunku całkowego dla całki Lebesgue'a mamy
$h(t)=\int_0^t h'(r)\,\dd r$. Odwrotnie każdej
$f\in L^2([0,T])$ odpowiada funkcja
$h_f(t)=\int_0^t f(r)\,\dd r$, która jest absolutnie
ciągła, zaczyna się od zera i ma pochodną $f$ prawie
wszędzie. Otrzymujemy izometryczną bijekcję z zupełną
przestrzenią $L^2$, a więc zupełność $\calH_T$.
Nierówność wynika z Cauchy'ego--Schwarza:
$|\int_s^t h'(r)\,\dd r|\le\sqrt{|t-s|}\,\|h'\|_2$.
Wreszcie $k_t'(r)=\ind{(0,t)}(r)$ niemal wszędzie, zatem
$\langle h,k_t\rangle_{\calH_T}=\int_0^t h'(r)\,\dd r=h(t)$.
\end{proof}

Pojawienie się $\min(s,t)$ nie jest przypadkowe. To jednocześnie
kowariancja procesu Wienera $\E B_sB_t$ i jądro reprodukujące
przestrzeni funkcji, o które wolno przesuwać miarę Wienera
bez zmiany jej zbiorów miary zero. Ten fakt stanie się treścią
następnego twierdzenia, a nie założeniem.

\begin{lemat}[Aproksymacja w przestrzeni Camerona--Martina]
\label{lem:cm-approx}
Dla $h\in\calH_T$ i podziału dyadycznego
$t_j^{(n)}=jT2^{-n}$ niech $h_n$ będzie interpolacją liniową
wartości $h(t_j^{(n)})$. Wtedy $h_n\in\calH_T$,
$h_n'$ jest schodkowa oraz
$\|h_n'-h'\|_2\to0$ i $\|h_n-h\|_\infty\to0$.
\end{lemat}
\begin{proof}
Na przedziale $(t_j^{(n)},t_{j+1}^{(n)})$ pochodna $h_n$
jest średnią funkcji $h'$ na tym przedziale:
\[
 h_n'|_{(t_j^{(n)},t_{j+1}^{(n)})}
 =\frac{1}{t_{j+1}^{(n)}-t_j^{(n)}}
 \int_{t_j^{(n)}}^{t_{j+1}^{(n)}}h'(s)\,\dd s.
\]
Operator zastępowania funkcji przez jej średnie przedziałowe
jest kontrakcją w $L^2$ z nierówności Jensena.
Dla funkcji ciągłych średnie zbiegają jednostajnie,
a funkcje ciągłe są gęste w $L^2$. Stąd dla każdego
$f\in L^2$ przybliżenia przedziałowe zbiegają w $L^2$ do $f$.
Biorąc $f=h'$ dostajemy pierwszą zbieżność. Druga wynika
z $\sup_t|h_n(t)-h(t)|\le\sqrt T\|h_n'-h'\|_2$.
\end{proof}

\section{Twierdzenie Camerona--Martina: gęstość przesuniętej miary}

Dla $h\in C_0([0,T])$ oznaczmy przez $\tau_h(w)=w+h$
przesunięcie trajektorii, a przez
$\mu_T^h=\mu_T\circ\tau_h^{-1}$ prawo $B+h$.
Nie należy mylić $\mu_T^h$ z przesunięciem samego argumentu
w całce Itô. Teraz przesuwamy \emph{całą miarę}.

\begin{lemat}[Skończeniewymiarowa translacja gaussowska]
\label{lem:gauss-shift}
Niech $Y\sim\Normal(0,D)$ w $\R^n$, gdzie $D$ jest
symetryczną dodatnio określoną macierzą, i niech
$a\in\R^n$. Dla każdej ograniczonej borelowskiej
$F:\R^n\to\R$ mamy
\[
\E F(Y+a)=\E\left[F(Y)\exp\left(
 a^{\mathsf T}D^{-1}Y-\frac12a^{\mathsf T}D^{-1}a
\right)\right].
\]
\end{lemat}
\begin{proof}
Gęstość $Y+a$ w punkcie $y$ jest proporcjonalna do
$\exp[-\frac12(y-a)^{\mathsf T}D^{-1}(y-a)]$.
Iloraz tej gęstości przez gęstość $Y$ jest równy
$\exp(a^{\mathsf T}D^{-1}y-\frac12a^{\mathsf T}D^{-1}a)$.
Pomnożenie przez $F$ i scałkowanie kończy dowód.
\end{proof}

\begin{twierdzenie}[Cameron--Martin na przedziale skończonym]
\label{th:cameron-martin}
Niech $T>0$ i $h\in\calH_T$.
Miary $\mu_T^h$ i $\mu_T$ są równoważne, a
\[
\frac{\dd\mu_T^h}{\dd\mu_T}(w)
=\exp\left(I(h')(w)-\frac12\|h'\|_2^2\right)
\quad\text{dla }\mu_T\text{-prawie każdego }w.
\tag{9.1}\label{eq:cm-density}
\]
Innymi słowy, dla każdego ograniczonego borelowskiego
funkcjonału $F:C_0([0,T])\to\R$ zachodzi
\[
\int F(w+h)\,\dd\mu_T(w)
=\int F(w)\exp\left(I(h')(w)-\frac12\|h'\|_2^2\right)
\,\dd\mu_T(w).
\tag{9.2}\label{eq:cm-identity}
\]
\end{twierdzenie}
\begin{proof}
\emph{Krok 1: przesunięcia łamanymi.} Załóżmy najpierw, że
$h'$ jest funkcją schodkową stałą na przedziałach podziału
$0=t_0<\cdots<t_m=T$. Dla funkcjonału cylindrycznego
$F(w)=\varphi(w(s_1),\ldots,w(s_k))$ uzupełnijmy podział
o wszystkie $s_j$. Wektor niezależnych przyrostów $B$ na
nowym podziale ma diagonalną macierz kowariancji, której
przekątna składa się z długości przedziałów.
Przesunięcie przez $h$ dodaje do każdego przyrostu
$\Delta h_j=\int_{t_j}^{t_{j+1}}h'(s)\,\dd s$.
Lemat \ref{lem:gauss-shift} zastosowany do tych przyrostów
daje gęstość
\[
\exp\left(\sum_j\frac{\Delta h_j}{\Delta t_j}\Delta B_j
 -\frac12\sum_j\frac{(\Delta h_j)^2}{\Delta t_j}\right)
=\exp\left(I(h')-\frac12\|h'\|_2^2\right).
\]
Tożsamość \eqref{eq:cm-identity} jest więc prawdziwa dla
wszystkich ograniczonych borelowskich funkcji cylindrycznych.

\emph{Krok 2: aproksymacja dowolnego $h\in\calH_T$.}
Niech $h_n$ będą interpolantami z lematu \ref{lem:cm-approx},
$f_n=h_n'$, a $f=h'$. Z izometrii Wienera
$I(f_n)\to I(f)$ w $L^2(\mu_T)$ i w prawdopodobieństwie,
a $\|f_n\|_2^2\to\|f\|_2^2$. Wykładniki
$Z_n=\exp(I(f_n)-\|f_n\|_2^2/2)$ zbiegają więc
w prawdopodobieństwie do
$Z=\exp(I(f)-\|f\|_2^2/2)$.
Ponadto dla zmiennej normalnej $I(f_n)$ obliczamy
$\E Z_n^2=e^{\|f_n\|_2^2}$; ta rodzina drugich
momentów jest ograniczona. Zatem $(Z_n)$ jest jednostajnie
całkowalna i $Z_n\to Z$ w $L^1(\mu_T)$.

Jeśli $F$ jest ograniczoną ciągłą funkcją cylindryczną,
to $F(B+h_n)\to F(B+h)$ punktowo, gdyż $h_n\to h$
jednostajnie. Zbieżność zdominowana daje zbieżność
lewych stron tożsamości dla $h_n$, a zbieżność $Z_n$
w $L^1$ daje zbieżność prawych stron.
Wobec tego \eqref{eq:cm-identity} zachodzi dla $h$ i
wszystkich ograniczonych ciągłych funkcji cylindrycznych.
Funkcje te rozdzielają miary na $C_0([0,T])$:
wyznaczają rozkłady skończeniewymiarowe, a te wyznaczają
miarę na $\mathcal B_T$ na mocy twierdzenia \ref{th:wiener-cyl}.
Zatem dwie miary po obu stronach \eqref{eq:cm-identity}
są równe, co rozszerza wzór na wszystkie ograniczone
borelowskie $F$.

\emph{Krok 3: równoważność.} Wykładnik $Z$ jest ściśle
dodatni i skończony $\mu_T$-prawie na pewno.
Ponieważ $\dd\mu_T^h=Z\,\dd\mu_T$, otrzymujemy
$\mu_T^h\ll\mu_T$; dodatniość $Z$ implikuje też
$\mu_T\ll\mu_T^h$.
\end{proof}

\begin{uwaga}[Znak w wykładniku i związek z Girsanowem]
W gęstości prawa $B+h$ względem prawa $B$ stoi
$+I(h')$. Jest to zgodne z rozdziałem \ref{ch:girsanov},
jeśli uwzględnimy kierunek, w którym zmieniamy miarę.
Prawo przesuniętej trajektorii pod starą miarą jest tym
samym, co odpowiednio przeważone prawo nieprzesuniętej
trajektorii pod nową miarą. Mylenie tych dwóch kierunków
prowadzi do błędów znaku.
\end{uwaga}

\section{Przesunięcia spoza przestrzeni Camerona--Martina}

\begin{twierdzenie}[Kryterium singularności przesunięcia]
\label{th:cm-singular}
Niech $h\in C_0([0,T])$. Jeśli $h\notin\calH_T$,
to miary $\mu_T^h$ i $\mu_T$ są wzajemnie singularne.
Dokładniej, istnieje zbiór borelowski $A\subset C_0([0,T])$
taki, że $\mu_T(A)=0$ i $\mu_T^h(A)=1$.
\end{twierdzenie}
\begin{proof}
Niech $t_j^{(n)}=jT2^{-n}$ i zdefiniujmy funkcję
schodkową $f_n$ jako nachylenia interpolantu $h$ na
przedziałach dyadycznych, bez zakładania istnienia $h'$.
Oznaczmy energię dyskretną
\[
E_n=\|f_n\|_2^2=
\sum_{j=0}^{2^n-1}\frac{[h(t_{j+1}^{(n)})-h(t_j^{(n)})]^2}
{t_{j+1}^{(n)}-t_j^{(n)}}.
\]
\emph{Krok 1: $E_n\to\infty$.} Każda średnia $f_n$ na
przedziale siatki jest średnią $f_{n+1}$ na dwóch
podprzedziałach. Z nierówności Jensena wynika
$E_n\le E_{n+1}$. Gdyby $(E_n)$ było ograniczone,
to $(f_n)$ byłoby ograniczone w przestrzeni Hilberta $L^2$.
Istniałby słabo zbieżny podciąg $f_{n_k}\rightharpoonup f$.
Dla każdej dyadycznej chwili $t$ oraz dostatecznie dużego $k$
mamy $\int_0^t f_{n_k}(r)\,\dd r=h(t)$.
Przechodząc do granicy słabej przy funkcji testowej
$\ind{[0,t]}\in L^2$ dostajemy $\int_0^t f=h(t)$.
Obie strony są ciągłe w $t$, a chwile dyadyczne są gęste,
więc równość zachodzi dla wszystkich $t$.
Oznaczałoby to $h\in\calH_T$, sprzecznie z założeniem.
Stąd $E_n\uparrow\infty$.

\emph{Krok 2: test rozdzielający prawa.} Niech
\[
G_n(w)=\sum_{j=0}^{2^n-1}\frac{\Delta_j h}{\Delta_j t}
\,[w(t_{j+1}^{(n)})-w(t_j^{(n)})].
\]
Jest to ciągła funkcja na $C_0([0,T])$.
Pod $\mu_T$ zmienna $G_n$ ma rozkład $\Normal(0,E_n)$,
ponieważ jest kombinacją niezależnych przyrostów Wienera.
Natomiast $G_n(w+h)=G_n(w)+E_n$, więc pod $\mu_T^h$
ma rozkład $\Normal(E_n,E_n)$.

\emph{Krok 3: Borel--Cantelli.} Wybierzmy rosnący podciąg
$n_k$ taki, że $E_{n_k}\ge k^3$.
Z nierówności Czebyszewa
\[
\mu_T\{G_{n_k}>E_{n_k}/2\}\le\frac4{E_{n_k}},\qquad
\mu_T^h\{G_{n_k}\le E_{n_k}/2\}\le\frac4{E_{n_k}}.
\]
Prawa strona jest sumowalna po $k$. Pierwszy lemat
Borela--Cantellego oznacza, że pod $\mu_T$ nierówność
$G_{n_k}>E_{n_k}/2$ zachodzi co najwyżej skończenie wiele razy,
a pod $\mu_T^h$ zachodzi od pewnego miejsca zawsze.
Zbiór
$A=\bigcup_{K\ge1}\bigcap_{k\ge K}\{G_{n_k}>E_{n_k}/2\}$
ma zatem $\mu_T(A)=0$ i $\mu_T^h(A)=1$.
\end{proof}

\begin{wniosek}[Pełna dychotomia przesunięć]
Dla każdego $h\in C_0([0,T])$ zachodzi dokładnie jedna
z dwóch sytuacji: $h\in\calH_T$ i wtedy $\mu_T^h\sim\mu_T$,
albo $h\notin\calH_T$ i wtedy $\mu_T^h\perp\mu_T$.
\end{wniosek}

\section{Przykłady rachunkowe i zadania rozwiązane}

\begin{zadanie}[Przesunięcie liniowe i zdarzenie końcowe]
Niech $h(t)=ct$, $c\in\R$, i $a\in\R$.
Wyznaczyć gęstość prawa $B+h$ oraz
$\mu_T^h\{w:w(T)>a\}$.

\emph{Rozwiązanie.} Ponieważ $h'=c$,
$\|h\|_{\calH_T}^2=c^2T$ i $I(h')=cB_T$.
Stąd $\dd\mu_T^h/\dd\mu_T=e^{cB_T-c^2T/2}$.
Dla zdarzenia końcowego możemy liczyć bezpośrednio:
$B_T+cT\sim\Normal(cT,T)$, a więc
\[
\mu_T^h\{w(T)>a\}=1-\Phi\left(\frac{a-cT}{\sqrt T}\right).
\]
Dla kontroli można też scałkować gęstość przeważoną:
$\int_a^\infty e^{cy-c^2T/2}p_T(y)\,\dd y$.
Po uzupełnieniu kwadratu wynik jest ten sam.
\end{zadanie}

\begin{zadanie}[Które funkcje potęgowe należą do przestrzeni Camerona--Martina?]
Niech $h(t)=t^\alpha$ na $[0,T]$, gdzie $\alpha>0$.
Wyznaczyć zakres $\alpha$, dla którego przesunięcie
miary Wienera przez $h$ zachowuje zbiory miary zero.

\emph{Rozwiązanie.} Dla $t>0$ mamy
$h'(t)=\alpha t^{\alpha-1}$. Funkcja $h$ jest
absolutnie ciągła dla $\alpha>0$; pozostaje sprawdzić
\[
\|h'\|_2^2=\alpha^2\int_0^T t^{2\alpha-2}\,\dd t.
\]
Całka jest skończona dokładnie wtedy, gdy
$2\alpha-2>-1$, czyli $\alpha>1/2$.
Dla $\alpha>1/2$ miary są równoważne, a dla
$0<\alpha\le1/2$ wzajemnie singularne.
Ciągłość $h$ nie wystarcza: kryterium dotyczy jej
\emph{energii kwadratowej}, nie samego modułu ciągłości.
\end{zadanie}

\begin{zadanie}[Maksimum ruchu Browna z dryfem]
Niech $a>0$, $c\in\R$ oraz
$M_T^c=\sup_{0\le s\le T}(B_s+cs)$.
Obliczyć $\Prob(M_T^c\ge a)$.

\emph{Rozwiązanie.} Twierdzenie \ref{th:cameron-martin}
z $h(t)=ct$ daje
\[
\Prob(M_T^c\ge a)=\E\bigl[
\ind{\{\sup_{s\le T}B_s\ge a\}}e^{cB_T-c^2T/2}\bigr].
\]
Na zdarzeniu $\{B_T\ge a\}$ całkujemy względem
$B_T\sim\Normal(0,T)$; na zdarzeniu
$\{\sup_{s\le T}B_s\ge a,\ B_T<a\}$
zasada odbicia daje gęstość końcową $p_T(2a-y)$
na $y<a$. Stąd
\[
\begin{split}
\Prob(M_T^c\ge a)
&=\int_a^\infty e^{cy-c^2T/2}p_T(y)\,\dd y\\
&\quad+\int_{-\infty}^{a}e^{cy-c^2T/2}p_T(2a-y)\,\dd y\\
&=\Phi\left(\frac{cT-a}{\sqrt T}\right)
+e^{2ca}\Phi\left(\frac{-a-cT}{\sqrt T}\right).
\end{split}
\]
Dla $c=0$ otrzymujemy klasyczny wynik
$2[1-\Phi(a/\sqrt T)]$.
\end{zadanie}

\begin{przyklad}[Entropia względna dwóch miar trajektorii]
Dla $h\in\calH_T$ iloraz gęstości opisuje także koszt
zmiany prawa trajektorii. Definiujemy entropię względną
$D(\mu_T^h\|\mu_T)=\int\log(\dd\mu_T^h/\dd\mu_T)
\,\dd\mu_T^h$. Z \eqref{eq:cm-density} i faktu, że
$I(h')$ pod przesuniętą miarą ma średnią $\|h'\|_2^2$,
dostajemy $D(\mu_T^h\|\mu_T)=\frac12\|h\|_{\calH_T}^2$.
Można ten ostatni fakt wykazać najpierw dla schodkowych
$h'$ przez sumy przyrostów, a następnie przejść do granicy
w $L^2$. Tym samym norma Camerona--Martina ma również
interpretację informacyjną.
\end{przyklad}

\begin{uwaga}[Most Browna a przesunięcie deterministyczne]
Prawo mostu Browna zakończonego w zerze w chwili $T$
jest singularne względem $\mu_T$, ponieważ pod pierwszym
prawem $w(T)=0$ prawie na pewno, a pod drugim
$\mu_T\{w(T)=0\}=0$. Nie ma sprzeczności z twierdzeniem
Camerona--Martina: most nie powstaje przez dodanie jednej
ustalonej funkcji $h$, lecz przez zmianę struktury kowariancji.
Na każdym krótszym przedziale $[0,t]$, $t<T$,
jego prawo jest jednak absolutnie ciągłe względem
prawa ruchu Browna, z gęstością
$p_{T-t}(-B_t)/p_T(0)$; wynika to z warunkowania
w końcowej chwili i wzoru na gęstości przejścia.
\end{uwaga}

\section{Rysunek: przestrzeń trajektorii i kierunki dopuszczalnych przesunięć}
\begin{center}
\begin{tikzpicture}[>=Latex,scale=0.94]
\draw[->] (0,0)--(9.0,0) node[right] {$t$};
\draw[->] (0,-1.6)--(0,2.4) node[above] {$w(t)$};
\draw[blue!70!black,very thick] plot[smooth] coordinates
{(0,0) (0.6,0.3) (1.0,-0.2) (1.4,0.45) (2.1,-0.10)
(2.8,0.7) (3.2,0.25) (4.0,0.9) (4.7,0.1) (5.3,0.65)
(6.1,0.4) (6.7,1.1) (7.3,0.35) (8.2,1.05)};
\draw[orange!85!black,very thick,dashed] plot[smooth] coordinates
{(0,0) (0.6,0.48) (1.0,0.10) (1.4,0.87) (2.1,0.53)
(2.8,1.54) (3.2,1.21) (4.0,2.1) (4.7,1.51) (5.3,2.18)
(6.1,2.14) (6.7,2.0) (7.3,1.93) (8.2,2.25)};
\node[blue!65!black,anchor=west] at (5.3,-1.1) {$w$};
\node[orange!70!black,anchor=west] at (5.3,-1.48) {$w+h$};
\end{tikzpicture}
\par\smallskip\small Schemat idei przesuwania trajektorii. Nie jest to losowa
symulacja; rysunek ilustruje różnicę między dodaniem funkcji $h$
a zmianą miary na przestrzeni wszystkich trajektorii.
\end{center}

\section{Podsumowanie i dalsza perspektywa}
Ustaliliśmy nie tylko \emph{jak} przeważać miarę Wienera,
ale także \emph{kiedy jest to możliwe bez zmiany zbiorów miary zero}.
Kryterium dokładnie wyznacza przestrzeń Camerona--Martina.
Nie jest to tożsamość z całką Itô: całka $I(h')$ pojawia się
we wzorze na pochodną Radona--Nikodyma \emph{dwóch miar}.
W ostatniej części tomu III wykorzystamy tę konstrukcję
w analizie całek po trajektoriach i półgrup Schrödingera.

% ===========================================================================
% ROZDZIAL 10 — Procesy Levy'ego
% ===========================================================================
